%%%%%%%%%%%%%%%%%%%%%%%%%%%%%%%%%%%%%%%%%%%%%%%%%%%%%%%
% NOTAS TEORIA DE TRANSPORTE OPTIMO
% FCEN-UBA, segundo cuatrimestre de 2026. Fuente unica de las notas
% y de las guias de ejercicios (herramientas/generar_practicas.py).
%%%%%%%%%%%%%%%%%%%%%%%%%%%%%%%%%%%%%%%%%%%%%%%%%%%%%%%

\documentclass[11pt,a4paper]{amsart}
\usepackage{xr-hyper}
\newcounter{chapter}
\usepackage{fancyhdr}
\usepackage{geometry} \geometry{top=2cm,bottom=2.5cm,left=2.5cm,right=2.5cm}

\usepackage{amssymb}
\usepackage{hyperref}
\usepackage{listings}
\usepackage{mathtools}
\usepackage{enumitem}

% Fuzz -------------------------------------------------------------------
% Line spacing -----------------------------------------------------------

%%%%%%%%%fin pegada%%%%%%%%%%%%%%%%%%5

\usepackage[T1]{fontenc}
\usepackage[spanish]{babel}
\usepackage{graphicx}
\usepackage{array}
\usepackage{tikz}

%% Colores del Instituto de Cálculo y del Departamento de Matemática (portada y enlaces)
\definecolor{icnavy}{HTML}{013274}
\definecolor{iccyan}{HTML}{029BDF}
\definecolor{icmag}{HTML}{E2017B}
\definecolor{dmred}{HTML}{A8000C}
\hypersetup{colorlinks=true,linkcolor=icnavy,citecolor=icnavy,urlcolor=icnavy}
\usetikzlibrary{babel}

%%%%%%%%%%%%%%%%%%%%%%%%%%%%%%%%%%%%%%%%%%%%%%%%%%%%%%%%%%%%%%%%%%%%%%%%%%
%% ESTILO COMÚN DE LAS FIGURAS TikZ
%% Todas las figuras se declaran con [fig] y no llevan scale=: el tamaño se
%% controla con las coordenadas, para que el texto de las etiquetas salga
%% siempre del mismo tamaño. La paleta deriva de los colores institucionales,
%% de modo que las figuras quedan en la misma gama que la portada y los
%% hipervínculos.
%%%%%%%%%%%%%%%%%%%%%%%%%%%%%%%%%%%%%%%%%%%%%%%%%%%%%%%%%%%%%%%%%%%%%%%%%%
\usetikzlibrary{arrows.meta,positioning,calc,patterns}

\colorlet{figmain}{icnavy}        % objeto principal (gráficos de funciones)
\colorlet{figacc}{icmag}          % objeto destacado (recta soporte, soporte del plan)
\colorlet{figalt}{iccyan}         % objeto secundario (pseudoinversa, segundo panel)
\colorlet{figwarn}{dmred}         % contraejemplos, lo que no puede pasar
\colorlet{figgrey}{black!55}      % anotaciones, construcciones auxiliares
\colorlet{figbg}{icnavy!6}        % fondo de las regiones X x Y
\colorlet{figfill}{iccyan!15}     % relleno de regiones destacadas

\tikzset{
  fig/.style={
    scale=1,
    line cap=round, line join=round,
    every node/.append style={font=\footnotesize},
    >={Stealth[length=5pt,width=4pt]},
  },
  eje/.style      ={->, thin, black!75},
  rotulo/.style   ={font=\footnotesize},
  anota/.style    ={font=\scriptsize, figgrey},
  graf/.style     ={very thick, figmain},
  grafalt/.style  ={very thick, figalt},
  soporte/.style  ={thick, figacc},
  destaca/.style  ={line width=3pt, figacc, line cap=round},
  aux/.style      ={dashed, thin, figgrey},
  grilla/.style   ={very thin, black!18},
  region/.style   ={fill=figbg},
  regionA/.style  ={fill=figfill},
  regionB/.style  ={fill=figwarn!12},
  prohibido/.style={thick, dotted, figwarn},
  pt/.style       ={circle, fill, inner sep=0pt, minimum size=3.6pt},
  ptvacio/.style  ={circle, draw, thick, fill=white, inner sep=0pt, minimum size=3.6pt},
  ptmasa/.style   ={circle, fill=figacc, inner sep=0pt},
}
\newcommand{\figpanelsep}{6cm}   % separación única entre paneles de una figura

%% En el índice de figuras, amsbook reserva 1.5pc para el número, que no alcanza
%% para «P.1» o «8.4» y pega el número contra el título. Lo ensanchamos.
\makeatletter
\renewcommand{\l@figure}{\@tocline{0}{3pt plus2pt}{0pt}{3pc}{}}
\makeatother



\parskip 3pt

%%
%% OPERADORES
%%
\renewcommand{\Re}{\operatorname{Re}}
\newcommand{\tr}{\operatorname{tr}}
\newcommand{\id}{\operatorname{id}}
\def\div{\operatorname {\mathrm{div}}}
\def\Dom{\operatorname {\mathrm{Dom}}}
\def\argmin{\operatorname {\mathrm{arg\, min}}}
\def\argmax{\operatorname {\mathrm{arg\, max}}}
\def\gen{\operatorname {\mathrm{gen}}}
\def\Im{\operatorname {\mathrm{Im}}}
\def\Nu{\operatorname {\mathrm{Nu}}}
\def\sop{\operatorname {\mathrm{sop}}}
\def\diam{\operatorname {\mathrm{diam}}}
\def\dist{\operatorname {\mathrm{dist}}}
\def\var{\operatorname {\mathrm{Var}}}
\def\codim{\operatorname {\mathrm{codim}}}
\def\pint{\operatorname {--\!\!\!\!\!\int\!\!\!\!\!--}}

%%
%% EPSILON
%%
\newcommand{\ve}{\varepsilon}

%%
%% LETRAS MATHBB
%%
\newcommand{\R}{{\mathbb R}}
\newcommand{\Z}{{\mathbb Z}}
\newcommand{\C}{{\mathbb C}}
\newcommand{\Q}{{\mathbb Q}}
\newcommand{\N}{{\mathbb N}}
\newcommand{\E}{{\mathbb E}}
\newcommand{\V}{{\mathbb V}}
\renewcommand{\P}{{\mathbb P}}
\newcommand{\II}{{\mathbb I}}
\newcommand{\OO}{{\mathbb O}}
\newcommand{\Sb}{{\mathbb S}}
\newcommand{\Tb}{{\mathbb T}}

%%
%% LETRAS CALIFGRAFICAS
%%
\newcommand{\F}{{\mathcal F}}
\newcommand{\Hc}{{\mathcal H}}
\newcommand{\Tc}{{\mathcal T}}
\newcommand{\Ec}{{\mathcal E}}
\newcommand{\Vc}{{\mathcal V}}
\newcommand{\A}{{\mathcal A}}
\newcommand{\J}{{\mathcal J}}
\newcommand{\G}{{\mathcal G}}
\newcommand{\U}{{\mathcal U}}
\newcommand{\B}{{\mathcal B}}
\newcommand{\Sw}{{\mathcal S}}
\newcommand{\D}{{\mathcal D}}
\newcommand{\M}{{\mathcal M}}
\newcommand{\LL}{{\mathcal L}}

%%
%% LETRAS MATHBF
%%
\newcommand{\ab}{{\mathbf a}}
\newcommand{\bb}{{\mathbf b}}
\newcommand{\cc}{{\mathbf c}}
\newcommand{\ee}{{\mathbf e}}
\newcommand{\ff}{{\mathbf f}}
\newcommand{\gb}{{\mathbf g}}
\newcommand{\hh}{{\mathbf h}}
\newcommand{\n}{{\mathbf n}}
\newcommand{\mm}{{\mathbf m}}
\newcommand{\pp}{{\mathbf p}}
\newcommand{\qq}{{\mathbf q}}
\newcommand{\sbf}{{\mathbf s}}
\newcommand{\tb}{{\mathbf t}}
\newcommand{\uu}{{\mathbf u}}
\newcommand{\vv}{{\mathbf v}}
\newcommand{\ww}{{\mathbf w}}
\newcommand{\xx}{{\mathbf x}}
\newcommand{\yy}{{\mathbf y}}
\newcommand{\zz}{{\mathbf z}}
\newcommand{\I}{{\mathbf 1}}
\newcommand{\Pb}{{\mathbf P}}
\newcommand{\XX}{{\mathbf X}}
\newcommand{\YY}{{\mathbf Y}}

%%
%% LETRAS MATHFRAK
%%

\newcommand{\Gk}{{\mathfrak G}}


%%
%% TEOREMAS
%%
\newtheorem{teo}{Teorema}[section]
\newtheorem{lema}[teo]{Lema}
\newtheorem{prop}[teo]{Proposición}
\newtheorem{cor}[teo]{Corolario}

%%
%% REMARK
%%
\theoremstyle{remark}
\newtheorem{obs}[teo]{Observación}

%%
%% DEFINICION
%%
\theoremstyle{definition}
\newtheorem{defi}[teo]{Definición}
\newtheorem{ejem}[teo]{Ejemplo}
\newtheorem{ejer}[teo]{Ejercicio}
\newtheorem{nota}[teo]{Notación}
\newtheorem{ejercicio}{Ejercicio}[chapter]


%%
%% NUMERACION
%%
\numberwithin{subsection}{section}
\numberwithin{equation}{section}

%%
%% TITULO
%%


\vfuzz7pt \hfuzz7pt
\externaldocument[N-]{../notas/Notas-TO}
\newcommand{\encabezado}[1]{%
\noindent \includegraphics[scale=2.1]{logo-dm-wide.png} \hskip 9.5cm
\includegraphics[scale=0.08]{logo-exactas-uba.jpg}
\smallskip
\usefont{T1}{phv}{m}{n}
\begin{center}{\Large\sc Teoría de Transporte Óptimo}\end{center}
\bigskip
\noindent{\sc #1}\\[-0.4cm]
\hrule
\normalfont}
\newcommand{\cuadernolink}[2]{\noindent\textbf{Cuaderno:} \emph{#1} (\href{https://colab.research.google.com/github/jfbonder/transporte-optimo/blob/main/cuadernos/#2}{abrir en Colab}).\par\medskip}

\begin{document}
\renewcommand{\thechapter}{2}
\encabezado{Práctica 2: Existencia de planes óptimos}
\bigskip
\pagestyle{fancy}
\fancyhead[L]{{\footnotesize{\usefont{T1}{phv}{m}{n}{\sc Teoría de Transporte Óptimo}}}}
\fancyhead[R]{\footnotesize \usefont{T1}{phv}{m}{n}{\sc Práctica 2}}
\thispagestyle{plain}
\noindent{\footnotesize Los ejercicios son los de la sección de ejercicios del Capítulo~2 de las notas del curso, con la misma numeración; las referencias a teoremas, proposiciones y secciones remiten a las notas (\url{https://github.com/jfbonder/transporte-optimo}). ($\bigstar$) ejercicio de mayor dificultad o que excede el programa.}
\medskip

\cuadernolink{Monge y Kantorovich}{01-monge-kantorovich.ipynb}


\begin{ejercicio}[Convergencia débil: ejemplos]\label{ejer.debil.ejemplos}
Decidir si las siguientes sucesiones convergen débilmente en $\mathcal P(\R)$
y, en caso afirmativo, hallar el límite:
\[
\begin{aligned}
&\text{(a) }\delta_{1/n};\quad
\text{(b) }\delta_n;\quad
\text{(c) }\tfrac12\delta_0+\tfrac12\delta_{1/n};\\[2pt]
&\text{(d) }\tfrac{n}{n-1}\LL^1\lfloor_{[0,1-1/n]};\quad
\text{(e) }n\,\LL^1\lfloor_{[0,1/n]};\quad
\text{(f) }\tfrac1n\sum_{k=1}^n\delta_{k/n}.
\end{aligned}
\]
En (b), probar que ninguna subsucesión converge débilmente a una medida de
probabilidad, y relacionarlo con la falta de ajuste (\emph{tightness}).
\end{ejercicio}

\begin{ejercicio}[Continuidad del push-forward]\label{ejer.pushforward.continuo}
Sean $X$, $Y$ espacios métricos y $T\colon X\to Y$ continua.
\begin{enumerate}[label=(\alph*)]
\item Probar que si $\mu_n\rightharpoonup\mu$ en $\mathcal P(X)$ entonces
$T_\#\mu_n\rightharpoonup T_\#\mu$ en $\mathcal P(Y)$.
\item Concluir que si $\mu_m=\frac1m\sum_{j=1}^m\delta_{x_j}$ converge
débilmente a $\mu$ (por ejemplo, si los $x_j$ son muestras independientes de
$\mu$, por la ley de los grandes números), entonces
$\frac1m\sum_j\delta_{T(x_j)}\rightharpoonup T_\#\mu$: «los mapas de
transporte mandan muestras en muestras».
\item Mostrar con un ejemplo que (a) falla si $T$ es solo medible.
\end{enumerate}
\end{ejercicio}

\begin{ejercicio}[Los mapas no son compactos]\label{ejer.mapas.no.compactos}
\leavevmode
\begin{enumerate}[label=(\alph*)]
\item Construir una sucesión $f_n\colon[0,1]\to[0,1]$ tal que
$(f_n)_\#\LL^1\lfloor_{[0,1]}=\LL^1\lfloor_{[0,1]}$ para todo $n$ y
$f_n\rightharpoonup\frac12$ débilmente en $L^2(0,1)$, es decir,
$\int f_n g\to\frac12\int g$ para toda $g\in L^2(0,1)$.
\emph{Sugerencia:} funciones «diente de sierra» de período $1/n$.
\item Sea $\pi_n=(\id,f_n)_\#\LL^1\lfloor_{[0,1]}$. Probar que
$\pi_n\rightharpoonup\LL^1\lfloor_{[0,1]}\otimes\LL^1\lfloor_{[0,1]}$, que no
está inducido por ningún mapa. Concluir que el conjunto de planes inducidos
por mapas no es débilmente cerrado, y comentar por qué esto obstruye el
método directo para el problema de Monge.
\item Probar que si $f\colon[0,1]\to[0,1]$ es de clase $C^1$ y
$f_\#\LL^1\lfloor_{[0,1]}=\LL^1\lfloor_{[0,1]}$, entonces $f(x)=x$ o
$f(x)=1-x$. Deducir que en (a) no es posible elegir las $f_n$ de clase $C^1$
(aunque sí lipschitzianas).
\end{enumerate}
\end{ejercicio}

\begin{ejercicio}[La variación total no sirve]\label{ejer.variacion.total}
La distancia de variación total en $\mathcal P(Z)$ es
$\|\alpha-\beta\|_{TV}=\sup_{A}|\alpha(A)-\beta(A)|$. Sea
$\mu=\LL^1\lfloor_{[0,1]}$ y $\pi_n=(\id,T_n)_\#\mu$ con
$T_n(x)=x+\frac1n \pmod 1$. Probar que $\pi_n\in\Pi(\mu,\mu)$, que
$\|\pi_n-\pi_m\|_{TV}=1$ para $n\ne m$, y concluir que $\Pi(\mu,\mu)$ no es
compacto para la variación total, aunque sí lo es para la topología débil
(Corolario~\ref{N-cor.Pi.compacto}). ¿A qué converge débilmente $\pi_n$?
\end{ejercicio}

\begin{ejercicio}[Regularización de Moreau--Yosida]\label{ejer.moreau.yosida}
Sea $(X,d)$ un espacio métrico y $g\colon X\to\R\cup\{+\infty\}$ propia
(no idénticamente $+\infty$) y acotada inferiormente. Para $k\ge0$ se define
\[
g_k(x)=\inf_{y\in X}\bigl\{g(y)+k\,d(x,y)\bigr\}.
\]
\begin{enumerate}[label=(\alph*)]
\item Probar que $g_k$ es finita, $k$-Lipschitz, $g_k\le g$ y $g_k\ge\inf g$.
\item Si además $g$ es semicontinua inferiormente, probar que $g_k(x)\nearrow g(x)$
para todo $x$ cuando $k\to\infty$.
\item Deducir que toda función semicontinua inferiormente y acotada
inferiormente en un espacio métrico es límite puntual creciente de una
sucesión de funciones Lipschitz y acotadas ($\min\{g_k,k\}$).
\end{enumerate}
\end{ejercicio}

\begin{ejercicio}[Integrales de funciones semicontinuas]\label{ejer.sci.integral}
Sea $X$ un espacio métrico, $f\colon X\to\R\cup\{+\infty\}$ semicontinua
inferiormente y acotada inferiormente, y $\mu\in\mathcal P(X)$.
\begin{enumerate}[label=(\alph*)]
\item Probar que
$\int_Xf\,d\mu=\sup\bigl\{\int_Xg\,d\mu:\ g\in C_b(X),\ g\le f\bigr\}$.
\emph{Sugerencia:} Ejercicio~\ref{ejer.moreau.yosida} y convergencia
monótona.
\item Deducir de (a) el Teorema~\ref{N-teo.K.sci}: si $\pi_n\rightharpoonup\pi$
en $\mathcal P(X\times Y)$ y $c$ es semicontinua inferiormente y acotada
inferiormente, entonces $\int c\,d\pi\le\liminf_n\int c\,d\pi_n$.
\item Mostrar con un ejemplo que la semicontinuidad inferior del funcional
de Kantorovich falla si $c$ es solo medible y acotada.
\end{enumerate}
\end{ejercicio}

\begin{ejercicio}[Estabilidad de los planes óptimos]\label{ejer.estabilidad.costos}
Sean $X$, $Y$ compactos, $\mu\in\mathcal P(X)$, $\nu\in\mathcal P(Y)$, y
$c_n,c\colon X\times Y\to\R$ continuas con $c_n\to c$ uniformemente. Sea
$\gamma_n$ un plan óptimo para $c_n$.
\begin{enumerate}[label=(\alph*)]
\item Probar que toda subsucesión de $(\gamma_n)$ tiene una subsucesión que
converge débilmente a un plan $\gamma$ óptimo para $c$, y que
$\min_{\Pi(\mu,\nu)}\int c_n\,d\gamma\to\min_{\Pi(\mu,\nu)}\int c\,d\gamma$.
\item Aplicarlo a $c_\ve(x,y)=|x-y|^{1+\ve}$ y a
$c_\ve(x,y)=\sqrt{\ve^2+|x-y|^2}$ con $\ve\to0$: los planes óptimos
convergen, salvo subsucesión, a planes óptimos para $|x-y|$. Como para
$|x-y|$ el plan óptimo no es único en general, estos límites
\emph{seleccionan} planes óptimos particulares; ¿qué se puede decir del
límite $\ve\to\infty$ en el segundo caso, una vez normalizado el costo?
\item ($\bigstar$) Si $\mu$ y $\nu$ tienen soportes disjuntos y
$\mu\ll\LL^d$, el límite en (b) es único y es el plan «monótono a lo largo
de rayos»; ver \cite[Sec.~3.1]{Santambrogio}.
\end{enumerate}
\end{ejercicio}

\begin{ejercicio}[$\bigstar$ La convergencia débil es metrizable]\label{ejer.debil.metrizable}
Sea $X$ un espacio métrico completo y separable. Probar que existe una
métrica $d$ en $\mathcal P(X)$ tal que $\mu_n\rightharpoonup\mu$ si y solo
si $d(\mu_n,\mu)\to0$. \emph{Sugerencia:} tomar una sucesión
$(g_k)\subset C_b(X)$ densa en un sentido adecuado (por ejemplo, densa en
$\operatorname{Lip}_b(X)$ para la convergencia uniforme sobre compactos) y
definir $d(\mu,\nu)=\sum_k2^{-k}\min\{1,|\int g_k\,d\mu-\int g_k\,d\nu|\}$;
ver \cite[Sec.~5.1]{AGS2008}. Compárese con la métrica $W_p$ del
Capítulo~\ref{N-cap.Wp}, que metriza la convergencia débil junto con la de
los momentos.
\end{ejercicio}

\subsection*{Ejercicios computacionales}

Los ejercicios siguientes se resuelven con el cuaderno \emph{Monge y
Kantorovich}, que acompaña los Capítulos~\ref{N-cap.intro} y~\ref{N-cap.existencia}
(véase el Apéndice~\ref{N-ap.cuadernos}); los enunciados figuran también al
final del cuaderno. Las referencias a secciones son a las del cuaderno.

\begin{ejercicio}[Asignación y Kantorovich]\label{ejc.MK.asignacion}
Para $n$ puntos con masas uniformes, verificar que el método exacto
(\texttt{ot.emd}) devuelve siempre una matriz de permutación dividida por $n$
y que su costo coincide con el del problema de asignación resuelto con la
rutina de SciPy, como predice el teorema de Birkhoff
(Ejercicio~\ref{N-ejer.biestocasticas}). Para masas genéricas, contar las
entradas no nulas del plan óptimo y comprobar la cota $n+m-1$. Construir un
ejemplo con $n=m=3$ en el que el plan óptimo tenga exactamente $5$ entradas no
nulas, y otro en el que tenga menos.
\end{ejercicio}

\begin{ejercicio}[El segmento $\Pi(\mu,\nu)$]\label{ejc.MK.segmento}
Para $\mu=\nu=\frac12(\delta_0+\delta_1)$ y los costos $c=|x-y|^2$,
$c=\mathbf 1_{\{x\ne y\}}$ y $c=-|x-y|^2$, graficar el costo como función del
parámetro $t$ que recorre el segmento $\Pi(\mu,\nu)$ e identificar todos los
planes óptimos en cada caso. Comprobar que en los tres el costo es afín en
$t$ y no constante, de modo que el mínimo se alcanza en un extremo del
segmento. ¿Qué condición sobre la matriz $(c_{ij})$ hace que \emph{todos} los
planes sean óptimos? Exhibir un costo que la cumpla y que no sea constante.
\end{ejercicio}

\begin{ejercicio}[Birkhoff]\label{ejc.MK.birkhoff}
Generar una matriz biestocástica $M$ como combinación convexa aleatoria de
cinco matrices de permutación. Para un costo aleatorio $C$, comprobar que el
plan óptimo del programa lineal es una permutación y que $\langle C,M\rangle$
es estrictamente mayor que el óptimo (salvo coincidencias). Implementar la
descomposición de Birkhoff: encontrar una permutación $\sigma$ con
$M_{i\sigma(i)}>0$ para todo $i$ (resolver un problema de asignación con costo
$-\log M_{ij}$, reemplazando los ceros por un valor finito grande, por
ejemplo con
\texttt{np.where(M > 0, -np.log(np.maximum(M, 1e-300)), 1e6)}: así el óptimo
evita las entradas nulas siempre que exista una permutación admisible, cosa
que el teorema de Birkhoff garantiza; con \texttt{inf} en cambio
\texttt{scipy.optimize.linear\_sum\_assignment} falla), restar $t\,M(\sigma)$ con $t=\min_iM_{i\sigma(i)}$ y repetir hasta
llegar a cero. ¿Cuántas permutaciones aparecen?
\end{ejercicio}

\begin{ejercicio}[Costos cóncavos]\label{ejc.MK.concavos}
Con las dos densidades superpuestas de la Sección~5 del cuaderno, calcular la
masa que el plan óptimo de $|x-y|^p$ deja en la celda diagonal de la
discretización como función de $p$ y graficarla, para $p\in(0,1)\cup(1,2]$,
marcando $p=1$ aparte: allí hay muchos planes óptimos y el valor que devuelve
el método exacto no está determinado. Repetir con dos tamaños de grilla
($N=100$ y $N=200$) para ver qué valores son estables y cuáles no; comparar
con la masa común $\sum_i\min(a_i,b_i)$. Explicar el salto en $p=1$ usando la
Proposición~\ref{N-prop:convexo-no-estricto} y el
Ejercicio~\ref{N-ejer.1D.concavo}.
\end{ejercicio}

\begin{ejercicio}[Convergencia débil de planes]\label{ejc.MK.debil}
Con las funciones «diente de sierra» $f_n$ del
Ejercicio~\ref{ejer.mapas.no.compactos}, verificar numéricamente que
$\int_0^1\phi(x)\psi(f_n(x))\,dx\to\int_0^1\phi\int_0^1\psi$ para varias
$\phi,\psi$ continuas (por ejemplo, polinomios y senos). Repetir con
$f_n(x)=x+\frac1n \bmod 1$: ¿a qué plan convergen ahora los
$\pi_n=(\id,f_n)_\#\LL^1\lfloor_{[0,1]}$? ¿Está inducido por un mapa?
\end{ejercicio}

\bibliographystyle{plain}
\bibliography{../notas/biblio}

\end{document}
